\documentclass[10pt]{article}
\usepackage[T1]{fontenc}
\usepackage[utf8]{inputenc}
\usepackage[french]{babel}
\usepackage{amsmath,amssymb,mathrsfs,mathpazo}
\usepackage{geometry}
\geometry{a4paper,margin=2cm}
\usepackage{graphicx,array,tabularx,tabularray,colortbl,multirow,multicol,booktabs,enumitem,pifont,ulem,fancybox,diagbox,cancel}
\usepackage{pstricks,pstricks-add,pst-plot,pst-tree,pst-node,pst-text,pst-func,pst-math,pst-bspline,pst-3dplot,pst-eucl,pst-solides3d}
\let\pscalculate\undefined
\usepackage{pst-bezier}
\usepackage{tikz,pgfplots,tkz-tab}
\usetikzlibrary{arrows,arrows.meta,patterns,calc,positioning,trees,decorations.markings,shapes.geometric,shapes.misc,shapes.arrows,intersections,decorations.pathmorphing,decorations.pathreplacing,backgrounds,shadows,babel}
\pgfplotsset{compat=1.18}
\usepgfplotslibrary{fillbetween}
\usepackage[np]{numprint}
\usepackage{listings,makecell,pas-tableur,twoopt}
\let\barre\undefined
\usepackage[pstricks]{tablvar}
\definecolor{darkgreen}{RGB}{0,100,0}
\providecommand{\red}{\color{red}}
\providecommand{\blue}{\color{blue}}
\providecommand{\green}{\color{green}}
\providecommand{\black}{\color{black}}
\providecommand{\white}{\color{white}}
\providecommand{\gray}{\color{gray}}
\providecommand{\cyan}{\color{cyan}}
\providecommand{\magenta}{\color{magenta}}
\providecommand{\yellow}{\color{yellow}}
\providecommand{\violet}{\color{violet}}
\providecommand{\orange}{\color{orange}}
\providecommand{\euro}{\texteuro}
\providecommand{\EUR}{\texteuro}
\providecommand{\ieme}{\textsuperscript{e}}
\providecommand{\degres}{\ensuremath{^\circ}}
\providecommand{\degree}{\ensuremath{^\circ}}
\providecommand{\vv}[1]{\overrightarrow{#1}}
\providecommand{\Coord}[1]{#1}
\providecommand{\decosix}{}
\providecommand{\starredbullet}{\ensuremath{\bullet}}
\pagestyle{empty}
\setlength{\parindent}{0pt}
\usepackage{adjustbox}\DeclareUnicodeCharacter{2236}{:}\DeclareUnicodeCharacter{2190}{\ensuremath{\leftarrow}}
\providecommand{\R}{\mathbb{R}}
\providecommand{\N}{\mathbb{N}}
\providecommand{\Z}{\mathbb{Z}}
\providecommand{\Q}{\mathbb{Q}}
\providecommand{\D}{\mathbb{D}}
\providecommand{\C}{\mathbb{C}}
\providecommand{\e}{{\mathrm{e}}}
\providecommand{\vect}[1]{\overrightarrow{#1}}
\providecommand{\vectt}[1]{\overrightarrow{\mathrm{#1}}}
\providecommand{\vvt}[1]{\overrightarrow{\mathrm{#1}}}
\providecommand{\np}[2][]{{#2~#1}}
\providecommand{\vv}[1]{{\overrightarrow{#1}}}
\renewcommand{\d}{{\mathrm{d}}}
\providecommand{\degres}{{^\circ}}
\providecommand{\degre}{{^\circ}}
\providecommand{\blue}{\color{blue}}
\providecommand{\red}{\color{red}}
\providecommand{\enskip}{\quad}
\providecommand{\llbracket}{\lbrack\!\lbrack}
\providecommand{\rrbracket}{\rbrack\!\rbrack}
\providecommand{\euro}{{\text{€}}}
\providecommand{\no}{{\text{n°}}}
\providecommand{\interval}[2]{{[#1;#2]}}
\providecommand{\intervFO}[2]{{[#1;#2[}}
\providecommand{\dotfill}{\ldots}

\providecommand{\ds}{\displaystyle}
\providecommand{\textding}[1]{#1}
\providecommand{\Oij}{\(\mathrm{O};\vec{\imath},\vec{\jmath}\)}
\providecommand{\Oijk}{\(\mathrm{O};\vec{\imath},\vec{\jmath},\vec{k}\)}
\providecommand{\Ouv}{\(\mathrm{O};\vec{u},\vec{v}\)}
\providecommand{\intervFF}[2]{\ensuremath{[#1;#2]}}
\renewcommand{\euro}{\ifmmode\text{\texteuro}\else\texteuro\fi}
\DeclareUnicodeCharacter{00B0}{\ifmmode^{\circ}\else\textdegree\fi}
\DeclareUnicodeCharacter{20AC}{\ifmmode\text{\texteuro}\else\texteuro\fi}
% Énoncés : Éducation nationale. Transcription et corrigés : APMEP, Denis Vergès et contributeurs. Découpage et classement : CoopMaths, Sébastien Lozano, Éric Elter et contributeurs. Adaptation de la présentation : Maths Explorer.
% CC BY-SA : https://coopmaths.fr/www/about/
% https://coopmaths.fr/alea/static/bac/2021/tex/bac_2021_09_metropole_10_cor.tex
\begin{document}
\noindent\textit{APMEP / CoopMaths. CC BY-SA. Présentation adaptée par Maths Explorer.}\par\medskip
\medskip

\begin{tabular}{|l|}\hline
Principaux domaines abordés :\\
Fonction logarithme, limites, dérivation.\\ \hline
\end{tabular}

\medskip

\textbf{Partie 1}

\medskip




\[f(x) = \dfrac{2\ln (x) - 1}{x}.\]

\begin{center}
\psset{unit=1.75cm}
\begin{pspicture*}(-0.6,-2.2)(5,1)
\psaxes[linewidth=1.25pt,labelFontSize=\scriptstyle]{->}(0,0)(0,-2.1)(5,1)
\psplot[plotpoints=2000,linewidth=1.25pt,linecolor=red]{0.1}{5}{x ln 2 mul 1 sub x div}
\end{pspicture*}
\end{center}

\medskip

\begin{enumerate}
\item 
Dans $]0~;~+\infty[$, \, $f(x) = 0 \iff \dfrac{2\ln (x) - 1}{x} = 0$, on a donc 

$2\ln (x) - 1 = 0 \iff 2\ln (x) = 1 \iff \ln (x) = \dfrac{1}{2} \iff x  = \text{e}^{\frac{1}{2}}$.

$S = \left\{\text{e}^{\frac{1}{2}}\right\}$.

\emph{Rem. } $\text{e}^{\frac{1}{2}} = \sqrt{\text{e}} \approx 1,649 \approx 1,65$ au centième près.


\item 
$\bullet~~$Sur $\left]0~;~\text{e}^{\frac{1}{2}}\right[$, on a $f(x) < 0$ ;

$\bullet~~$Sur $\left]\text{e}^{\frac{1}{2}}~;~+ \infty\right[$, on a $f(x) > 0$ ;

$\bullet~~$ $f\left(\text{e}^{\frac{1}{2}}\right) = 0$.
\end{enumerate}

\bigskip

\textbf{Partie II}

\medskip



\[g(x) = [\ln (x)]^2 - \ln (x).\]

\begin{enumerate}
\item 
	\begin{enumerate}
		\item 
On a $\displaystyle\lim_{x \to 0} \ln x = - \infty$, d'où $\displaystyle\lim_{x \to 0} (\ln x)^2 = + \infty$ et $\displaystyle\lim_{x \to 0} (-\ln x) = + \infty$, donc par somme de limites :
		
$\displaystyle\lim_{x \to 0} g(x) = + \infty$.
		\item 
$g(x) = \ln (x) [\ln (x) - 1]$. Comme 

$\displaystyle\lim_{x \to + \infty} \ln (x)  = + \infty$ et $\displaystyle\lim_{x \to + \infty} \ln (x)  - 1 = + \infty$, on obtient par produit :

$\displaystyle\lim_{x \to + \infty} g(x)  = + \infty$.
	\end{enumerate}
\item 
La fonction $g(x) = \ln (x) [\ln (x) - 1]$ est dérivable comme produit de deux fonctions dérivables sur $]0~;~+ \infty[$ et sur cet intervalle :

$g'(x) = \dfrac{1}{x} \times [\ln (x) - 1] + \ln (x) \times \dfrac{1}{x} = \dfrac{1}{x} \left[\ln (x) - 1 + \ln (x)\right] = \dfrac{1}{x}\times (2\ln (x) - 1) = \dfrac{2\ln (x) - 1}{x} = f(x)$.

\item 

Le signe de $f(x) = g'(x)$ a été trouvé à la question 2 de la partie I ; on a donc :

$\bullet~~$Sur $\left]0~;~\text{e}^{\frac{1}{2}}\right[$, on a $g'(x) < 0$ : la fonction $g$ est strictement décroissante sur cet intervalle

$\bullet~~$Sur $\left]\text{e}^{\frac{1}{2}}~;~+ \infty\right[$, on a $g'(x) > 0$ : la fonction $g$ est strictement croissante sur cet intervalle

$\bullet~~$ $g'\left(\text{e}^{\frac{1}{2}}\right) = 0$ : $g\left(\text{e}^{\frac{1}{2}}\right) = \dfrac{1}{4} - \dfrac{1}{2} = - \dfrac{1}{4}$ est le minimum de la fonction $g$ sur $]0~;~+\infty[$.


\begin{center}
\psset{unit=1cm}
\begin{pspicture}(7,3)
\psframe(7,3)\psline(0,2)(7,2)(7,3)\psline(0,2.5)(7,2.5)\psline(1,0)(1,3)
\uput[u](0.5,2.4){$x$}~\uput[u](1.15,2.4){$0$}~\uput[u](4,2.4){$\text{e}^{\frac{1}{2}}$}~ \uput[u](6.5,2.4){$+ \infty$}
\uput[u](0.5,1.9){$g'(x)$}\uput[u](2.5,1.9){$-$}\uput[u](4,1.9){$0$}\uput[u](5.5,1.9){$+$}
\psline{->}(1.5,1.5)(3.5,0.5)\psline{->}(4.5,0.5)(6.5,1.5)
\uput[d](1.5,2){$+ \infty$}\uput[u](4,0){$- \frac{1}{4}$}\uput[d](6.5,2){$+ \infty$}
\rput(0.5,1){$g$}
\end{pspicture}
\end{center}


\item 
Comme $- \frac{1}{4} = - 0,25$, le tableau de variations montre que l'équation $g(x) = m$, avec $m > -0,25$ a deux solutions, l'une sur l'intervalle $\left]0~;~\text{e}^{\frac{1}{2}}\right[$, l'autre sur $\left]\text{e}^{\frac{1}{2}}~;~+ \infty\right[$.
\item 
Dans $]0~;~+ \infty[$, \, $g(x) = 0 \iff \ln (x) [\ln (x) - 1] = 0 \iff \left\{\begin{array}{l c l}
\ln (x)&=&0\\
\ln (x) - 1&=&0
\end{array}\right. \iff $

$\left\{\begin{array}{l c l}
\ln (x)&=&0\\
\ln (x) &=&1
\end{array}\right. \iff \left\{\begin{array}{l c l}
x&=&1\\
x &=&\text{e}
\end{array}\right.$

$S = \{1~;~\text{e}\}$.
\end{enumerate}

\end{document}
